\documentclass[11pt,a4paper]{article}

% Compilation conseillée : pdflatex documentation_generale_modele_hitran.tex
\usepackage[utf8]{inputenc}
\usepackage[T1]{fontenc}
\usepackage[french]{babel}
\usepackage{lmodern}
\usepackage{microtype}
\usepackage{geometry}
\usepackage{amsmath,amssymb}
\usepackage{siunitx}
\usepackage[version=4]{mhchem}
\usepackage{booktabs}
\usepackage{array}
\usepackage{enumitem}
\usepackage{xcolor}
\usepackage{graphicx}
\usepackage{fancyhdr}
\usepackage{hyperref}

\geometry{top=2.35cm,bottom=2.35cm,left=2.45cm,right=2.45cm}
\sisetup{locale=FR,per-mode=symbol}
\DeclareSIUnit\ppm{ppm}
\DeclareSIUnit\atm{atm}
\DeclareSIUnit\molecule{molécule}

\definecolor{bleuCAPECL}{RGB}{17,67,104}
\definecolor{bleupale}{RGB}{234,243,250}
\hypersetup{
  colorlinks=true,
  linkcolor=bleuCAPECL,
  urlcolor=blue!60!black,
  pdftitle={Documentation générale du modèle spectral d'émissivité fondé sur HITRAN}
}

\pagestyle{fancy}
\fancyhf{}
\lhead{CAPECL -- Projet climat 2026}
\rhead{Modèle HITRAN}
\cfoot{\thepage}
\setlength{\headheight}{14pt}

\newcommand{\HITRAN}{\textsc{Hitran}}
\newcommand{\HAPI}{\textsc{Hapi}}
\newcommand{\code}[1]{\texttt{#1}}
\newcommand{\important}[1]{%
  \par\medskip
  \noindent\fcolorbox{bleuCAPECL}{bleupale}{%
    \parbox{0.94\linewidth}{#1}}%
  \par\medskip
}

\title{\vspace{2.2cm}
  {\color{bleuCAPECL}\LARGE\bfseries Modèle spectral d'émissivité\\[0.25em]
  atmosphérique fondé sur \HITRAN}\par\vspace{0.8cm}
  {\large Documentation générale du modèle}}
\author{Projet climat 2026\\Groupe C -- Geckoquins}
\date{Juin 2026}

\begin{document}

\maketitle
\thispagestyle{empty}
\vfill
\begin{center}
\begin{minipage}{0.78\textwidth}
\small
Cette documentation complète les programmes du dossier
\code{modèle2\_hitran\_emissivité}. Elle expose la construction physique du
modèle, le rôle des données spectroscopiques et les précautions nécessaires à
l'interprétation des résultats. Elle n'est pas une explication ligne par
ligne du code.
\end{minipage}
\end{center}
\vfill
\newpage

\tableofcontents
\newpage

\section{Pourquoi un second modèle ?}

Le premier modèle représente chaque bande d'absorption par quelques pics
lisses. Cette construction permet de vérifier rapidement les ordres de
grandeur, mais elle efface la structure spectrale réelle et ne fait pas varier
les sections efficaces avec la pression et la température.

Le second modèle conserve l'atmosphère discrétisée et la loi de
Beer--Lambert, mais remplace les pics calibrés manuellement par des spectres
calculés raie par raie à partir de
\href{https://hitran.org/}{HITRANonline}. L'interface Python
\href{https://hitran.org/hapi/}{HAPI} assure l'accès aux données et le calcul
spectral.

Le contexte physique général est celui du
\href{https://science.nasa.gov/earth/earth-observatory/climate-and-earths-energy-budget/}{bilan
énergétique terrestre décrit par la NASA} : la vapeur d'eau, le dioxyde de
carbone, le méthane, l'ozone et d'autres gaz absorbent certaines longueurs
d'onde du rayonnement thermique. Le présent modèle calcule l'absorption d'une
couche, pas le bilan radiatif complet de la planète.

\section{Grandeur calculée et domaine}

Le résultat principal est l'absorptivité spectrale d'une couche homogène :

\begin{equation}
  \alpha(\lambda,z)=1-\exp[-\tau_{\mathrm{tot}}(\lambda,z)].
\end{equation}

Dans les fichiers, cette grandeur est appelée \code{emissivite}. L'égalité
entre absorptivité et émissivité est une application de la loi de Kirchhoff en
équilibre thermodynamique local. Elle ne transforme pas le résultat en flux
émis : un flux demanderait la fonction de Planck et la propagation du
rayonnement dans toute la colonne.

Le tableau actuel couvre :

\begin{itemize}[itemsep=2pt]
  \item quatre absorbeurs : \ce{H2O}, \ce{CO2}, \ce{O3} et \ce{CH4} ;
  \item les longueurs d'onde de 4 à \SI{100}{\micro\metre} ;
  \item les altitudes de 0 à \SI{100}{\kilo\metre} ;
  \item des couches verticales de \SI{1}{\kilo\metre} ;
  \item un pas de sortie de \SI{0.1}{\micro\metre} et un pas interne de
        \SI{0.05}{\per\centi\metre}.
\end{itemize}

La chaîne de modélisation est

\begin{equation*}
\boxed{\text{paramètres de raies \HITRAN}}
\longrightarrow
\boxed{\sigma_g(\widetilde\nu,T,P)}
\longrightarrow
\boxed{n_g(z)}
\longrightarrow
\boxed{\tau_g(\lambda,z)}
\longrightarrow
\boxed{\alpha(\lambda,z)}.
\end{equation*}

\section{Atmosphère discrétisée}

\subsection{Hypothèses de couche}

Chaque couche est considérée plane, homogène et isotherme sur le trajet
radiatif. La pression, la température, les fractions molaires et les sections
efficaces calculées à ces conditions sont constantes sur $\Delta z$.

La propagation est verticale et sans diffusion. Les nuages, les aérosols, la
réfraction et la courbure terrestre sont absents. Les quatre gaz sont traités
comme des absorbeurs indépendants dont les épaisseurs optiques s'additionnent.

\subsection{Température et pression}

Le profil externe emploie une atmosphère standard simplifiée, affine par
morceaux. Dans une zone de gradient $\gamma$ constant,

\begin{align}
  T(z)&=T_0+\gamma(z-z_0),\\
  P(z)&=P_0\left(\frac{T(z)}{T_0}\right)^{-gM/(R\gamma)}
  \quad (\gamma\neq0),
\end{align}

et dans une zone isotherme,

\begin{equation}
  P(z)=P_0\exp\!\left[-\frac{gM(z-z_0)}{RT_0}\right].
\end{equation}

Cette construction s'inspire de la
\href{https://www.ngdc.noaa.gov/stp/space-weather/online-publications/miscellaneous/us-standard-atmosphere-1976/us-standard-atmosphere_st76-1562_noaa.pdf}{U.S.
Standard Atmosphere 1976}, mais elle n'en est pas une reproduction. La
stratopause est simplifiée et le gradient au-dessus de
\SI{85}{\kilo\metre} est propre au projet.

\subsection{Profils de gaz}

Les fractions molaires $x_g=C_g\times10^{-6}$ proviennent de profils
paramétriques :

\begin{table}[ht]
\centering
\small
\begin{tabular}{>{\bfseries}l p{0.70\linewidth}}
\toprule
Gaz & Profil de mélange utilisé \\
\midrule
\ce{CO2} & \SI{425}{\ppm} jusqu'à \SI{85}{\kilo\metre}, puis décroissance
exponentielle de hauteur caractéristique \SI{10}{\kilo\metre}. \\
\ce{CH4} & \SI{1.9}{\ppm} jusqu'à \SI{11}{\kilo\metre}, puis décroissance
exponentielle de hauteur caractéristique \SI{25}{\kilo\metre}. \\
\ce{O3} & Fond de \SI{0.1}{\ppm}; pic gaussien entre 15 et
\SI{35}{\kilo\metre}, centré à \SI{25}{\kilo\metre} et de maximum
\SI{8}{\ppm}. \\
\ce{H2O} & Pression de vapeur saturante approchée par la formule de Magnus et
humidité relative décroissante dans la troposphère; plancher de
\SI{4}{\ppm} au-dessus. \\
\bottomrule
\end{tabular}
\caption{Paramétrisation des profils verticaux.}
\end{table}

Ces profils ne sont pas des observations météorologiques instantanées. Les
valeurs de référence de \ce{CO2} et \ce{CH4} peuvent être comparées aux
séries du
\href{https://gml.noaa.gov/ccgg/}{Global Greenhouse Gas Reference Network de
la NOAA}. Leur fonction ici est de définir une atmosphère de test cohérente.

\section{Données spectroscopiques HITRAN}

\subsection{Contenu de la base}

HITRAN fournit, pour chaque transition moléculaire, la position de la raie,
son intensité de référence, l'énergie de l'état inférieur, les paramètres
d'élargissement et de déplacement collisionnels ainsi que les informations
isotopologiques. Les définitions et les unités sont rassemblées sur la page
\href{https://hitran.org/docs/definitions-and-units/}{HITRANonline --
Definitions and units}.

\begin{table}[ht]
\centering
\begin{tabular}{lccc}
\toprule
Molécule & Identifiant HITRAN & Isotopologue & Nom local \\
\midrule
\ce{H2O} & 1 & 1 & \code{H2O} \\
\ce{CO2} & 2 & 1 & \code{CO2} \\
\ce{O3}  & 3 & 1 & \code{O3} \\
\ce{CH4} & 6 & 1 & \code{CH4} \\
\bottomrule
\end{tabular}
\caption{Espèces spectroscopiques retenues. L'identifiant isotopologique 1
désigne l'isotopologue terrestre le plus abondant de la molécule.}
\end{table}

Seul l'isotopologue principal de chaque gaz est inclus. Le dossier
\code{donnees\_hitran} constitue un cache local des tables téléchargées. Sa
présence évite une nouvelle requête, mais ne certifie ni l'édition de la base,
ni la date d'acquisition, ni la couverture exacte de chaque fichier.

\subsection{Longueur d'onde et nombre d'onde}

HAPI travaille en nombre d'onde $\widetilde\nu$, en
\si{\per\centi\metre}. Pour $\lambda$ exprimée en mètres,

\begin{equation}
  \widetilde\nu\;(\si{\per\centi\metre})=
  \frac{1}{100\,\lambda\;(\si{\metre})}.
\end{equation}

Ainsi, l'intervalle de 4 à \SI{100}{\micro\metre} correspond à 2500--
\SI{100}{\per\centi\metre}. Le calcul interne est effectué sur une grille
uniforme en nombre d'onde.

\subsection{Profil de Voigt}

Chaque raie est représentée par un profil de Voigt, convolution d'un profil
gaussien et d'un profil lorentzien. Schématiquement,

\begin{equation}
  \sigma_g(\widetilde\nu,T,P)=
  \sum_{j\in g}S_j(T)
  V\!\left(\widetilde\nu-\widetilde\nu_j;
  \gamma_{D,j}(T),\gamma_{L,j}(T,P)\right).
\end{equation}

Le terme gaussien représente l'élargissement Doppler dû au mouvement
thermique. Le terme lorentzien représente principalement l'élargissement par
collisions et dépend de la pression et du gaz perturbateur. Le détail des
paramètres est donné dans la documentation HITRAN et dans le
\href{https://hitran.org/static/hapi/hapi_manual.pdf}{manuel officiel de
HAPI}.

Avec l'option d'unités HITRAN, la sortie est une section efficace en
\si{\centi\metre\squared\per\molecule}. Le modèle applique

\begin{equation}
  \sigma_g\;(\si{\metre\squared\per\molecule})=
  10^{-4}\sigma_g\;(\si{\centi\metre\squared\per\molecule}).
\end{equation}

La pression locale est convertie en atmosphères avant le calcul :

\begin{equation}
  P_{\mathrm{atm}}=\frac{P_{\mathrm{Pa}}}{101325}.
\end{equation}

\section{De la section efficace à l'absorptivité}

La densité moléculaire totale de l'air est

\begin{equation}
  n_{\mathrm{air}}(z)=\frac{P(z)}{k_{\mathrm B}T(z)},
  \qquad
  k_{\mathrm B}=\SI{1.380649e-23}{\joule\per\kelvin},
\end{equation}

où la valeur exacte de $k_{\mathrm B}$ est documentée par le
\href{https://www.nist.gov/si-redefinition/kelvin/kelvin-boltzmann-constant}{NIST}.
Pour le gaz $g$,

\begin{equation}
  n_g(z)=x_g(z)n_{\mathrm{air}}(z).
\end{equation}

L'épaisseur optique d'une couche homogène vaut

\begin{equation}
  \tau_g(\lambda,z)=
  \sigma_g(\lambda,T(z),P(z))\,n_g(z)\,\Delta z.
\end{equation}

Les contributions sont additionnées,

\begin{equation}
  \tau_{\mathrm{tot}}=
  \tau_{\ce{CO2}}+\tau_{\ce{H2O}}+\tau_{\ce{CH4}}+\tau_{\ce{O3}},
\end{equation}

puis Beer--Lambert donne

\begin{equation}
  \mathcal T_\lambda=\exp(-\tau_{\mathrm{tot}}),
  \qquad
  \boxed{\alpha_\lambda=1-\exp(-\tau_{\mathrm{tot}})}.
\end{equation}

L'épaisseur optique est sans dimension. La formulation HITRAN équivalente,
fondée sur la densité de colonne $u=\int n_g\,\mathrm dl$, est présentée dans
la page
\href{https://hitran.org/docs/definitions-and-units/}{Definitions and units}.

\section{Résolution spectrale et interpolation}

Le spectre HAPI est calculé par pas de
\SI{0.05}{\per\centi\metre}. Il est ensuite interpolé aux nombres d'onde
correspondant à une grille de sortie uniforme en longueur d'onde, par pas de
\SI{0.1}{\micro\metre}.

Une grille uniforme en $\lambda$ n'est pas uniforme en $\widetilde\nu$ :

\begin{equation}
  \left|\frac{\mathrm d\widetilde\nu}{\mathrm d\lambda}\right|
  \propto\frac{1}{\lambda^2}.
\end{equation}

La résolution effective du CSV varie donc avec la longueur d'onde. Une simple
interpolation ponctuelle peut manquer des raies plus étroites que le pas de
sortie. Pour un calcul de flux ou une comparaison instrumentale, il est
préférable d'intégrer la transmission sur chaque bande spectrale plutôt que
d'échantillonner uniquement son centre.

\section{Lecture du tableau final}

La table \code{table\_emissivite\_hitran\_4\_100um\_0\_100km.csv} contient
\num{97061} lignes de données : 101 altitudes et 961 longueurs d'onde. Une
ligne contient les coordonnées spectrales et verticales, la pression, la
température, les concentrations, les quatre sections efficaces, les quatre
épaisseurs optiques partielles et l'absorptivité appelée \code{emissivite}.

L'épaisseur optique totale n'est pas enregistrée directement, mais se
reconstruit par

\begin{equation}
  \tau_{\mathrm{tot}}=\sum_g\tau_g
  =-\ln(1-\alpha),\qquad \alpha<1.
\end{equation}

\important{\textbf{Table à régénérer avant usage quantitatif.} Le CSV présent
dans le dossier contient encore \SI{425}{\ppm} de \ce{CO2} à
\SI{100}{\kilo\metre}, tandis que le profil actuel impose une décroissance
au-dessus de \SI{85}{\kilo\metre} et donne environ \SI{94.83}{\ppm} à
\SI{100}{\kilo\metre}. Le fichier et le profil source ne correspondent donc
plus exactement.}

Une ligne décrit une couche isolée. Le rayonnement sortant d'une colonne exige
d'ordonner les couches, d'atténuer l'émission de chacune par les couches
supérieures et d'employer sa température dans la fonction de Planck.

\section{Apport par rapport au modèle empirique}

\begin{table}[ht]
\centering
\small
\begin{tabular}{p{0.24\linewidth}p{0.32\linewidth}p{0.34\linewidth}}
\toprule
Aspect & Modèle empirique & Modèle HITRAN \\
\midrule
Spectre & Quelques enveloppes exponentielles & Transitions moléculaires raie
par raie \\
Dépendance en $P,T$ & Absente & Incluse dans le calcul Voigt \\
Traçabilité & Paramètres internes non sourcés & Base et publications HITRAN \\
Coût numérique & Faible & Élevé \\
Usage naturel & Compréhension et tests & Calcul spectral plus réaliste \\
Limite commune & Absorptivité d'une couche, pas un flux radiatif complet &
Absorptivité d'une couche, pas un flux radiatif complet \\
\bottomrule
\end{tabular}
\caption{Positionnement des deux modèles.}
\end{table}

Les figures du dossier \code{comparaison\_modeles\_sections\_efficaces}
montrent que l'enveloppe empirique peut localiser une grande bande, mais ne
reproduit ni la multitude des raies, ni toutes les bandes faibles, ni leurs
ailes spectrales. Elles constituent une validation qualitative, pas une
validation radiométrique du modèle complet.

\section{Limites du modèle}

Le recours à HITRAN améliore fortement la représentation spectrale, mais ne
transforme pas le calcul en modèle climatique complet. Les limites
principales sont :

\begin{itemize}[itemsep=2pt]
  \item un seul isotopologue par gaz ;
  \item un profil de Voigt avec les choix de diluant et de coupure de HAPI ;
  \item l'absence des continua, notamment celui de la vapeur d'eau, et des
        absorptions induites par collision ;
  \item des profils atmosphériques idéalisés, surtout au-dessus de
        \SI{85}{\kilo\metre} ;
  \item l'absence de nuages, d'aérosols et de diffusion ;
  \item des couches planes et homogènes et un trajet uniquement vertical ;
  \item une table de sortie trop peu résolue pour conserver toutes les raies ;
  \item l'absence de fonction de Planck, de flux montant et descendant et de
        bilan énergétique ;
  \item un cache local qui ne mémorise pas explicitement l'édition HITRAN et
        la date de téléchargement.
\end{itemize}

Les améliorations prioritaires sont la régénération du CSV avec le profil
actuel, l'export de $\tau_{\mathrm{tot}}$, l'étude de convergence spectrale et
l'ajout d'un véritable transfert radiatif multicouche.

\section{Reproductibilité et correspondance avec les fichiers}

L'environnement Python est décrit par \code{requirements.txt}. Au premier
téléchargement, HAPI demande un accès à HITRANonline et peut nécessiter une
clé d'API. Pour rendre un résultat reproductible, il faut archiver avec le
CSV :

\begin{itemize}[itemsep=2pt]
  \item la version de HAPI ;
  \item l'édition ou la date d'accès à HITRAN ;
  \item la plage et le pas en nombre d'onde ;
  \item les isotopologues choisis ;
  \item la version du profil atmosphérique ;
  \item l'épaisseur et l'altitude représentative des couches.
\end{itemize}

La carte d'accès scientifique aux fichiers est :

\begin{table}[ht]
\centering
\small
\begin{tabular}{p{0.45\linewidth}p{0.47\linewidth}}
\toprule
Fichier & Objet scientifique associé \\
\midrule
\code{1\_section\_efficace\_CO2\_15um.py} & Prototype autour de la bande de
\ce{CO2} à \SI{15}{\micro\metre}. \\
\code{2\_sections\_efficaces\_gaz\_hitran.py} & Acquisition des raies et
construction des quatre spectres Voigt. \\
\code{code\_final\_emissivite\_avec\_sections\_hitran.py} & Couplage des
spectres aux profils verticaux et à Beer--Lambert. \\
\code{generer\_table\_emissivite\_hitran.py} & Échantillonnage altitude--
longueur d'onde et export du CSV. \\
\code{donnees\_hitran/} & Tables locales lues par HAPI. \\
\bottomrule
\end{tabular}
\end{table}

\section{Sources et ressources}

\begin{itemize}[leftmargin=*,itemsep=3pt]
  \item \href{https://science.nasa.gov/earth/earth-observatory/climate-and-earths-energy-budget/}{NASA Earth Observatory -- Climate and Earth's Energy Budget} ;
  \item \href{https://www.ngdc.noaa.gov/stp/space-weather/online-publications/miscellaneous/us-standard-atmosphere-1976/us-standard-atmosphere_st76-1562_noaa.pdf}{NASA/NOAA -- U.S. Standard Atmosphere 1976} ;
  \item \href{https://gml.noaa.gov/ccgg/}{NOAA Global Monitoring Laboratory -- Carbon Cycle Greenhouse Gases} ;
  \item \href{https://www.nist.gov/si-redefinition/kelvin/kelvin-boltzmann-constant}{NIST -- Boltzmann constant and the kelvin} ;
  \item \href{https://hitran.org/docs/}{HITRANonline -- documentation générale} ;
  \item \href{https://hitran.org/docs/definitions-and-units/}{HITRANonline -- Definitions and units} ;
  \item \href{https://hitran.org/hapi/}{HAPI -- page officielle} et
        \href{https://hitran.org/static/hapi/hapi_manual.pdf}{manuel utilisateur} ;
  \item I.~E. Gordon et al.,
        \href{https://doi.org/10.1016/j.jqsrt.2021.107949}{The HITRAN2020 Molecular Spectroscopic Database}, 2022 ;
  \item R.~V. Kochanov et al.,
        \href{https://doi.org/10.1016/j.jqsrt.2016.03.005}{HITRAN Application Programming Interface (HAPI)}, 2016 ;
  \item O.~A. Alduchov et R.~E. Eskridge,
        \href{https://doi.org/10.1175/1520-0450(1996)035\%3C0601:imfaos\%3E2.0.co;2}{Improved Magnus Form Approximation of Saturation Vapor Pressure}, 1996.
\end{itemize}

\end{document}
